% Resumo (ABNT NBR 6028: 150–500 palavras, parágrafo único) e Abstract.
\noindent
Os distúrbios osteomusculares relacionados ao trabalho atingem a maior parte dos
trabalhadores de escritório e estão associados a posturas inadequadas mantidas
por longos períodos. Este trabalho apresenta um sistema autônomo que detecta
postura sentada prejudicial e alerta o usuário em tempo real por vibração. O
sistema combina uma cadeira de escritório instrumentada com oito sensores de
força resistivos (FSR402), quatro no assento e quatro no encosto, e uma unidade
vestível com unidade de medição inercial (IMU), posicionada logo abaixo do
pescoço. Ambas usam o microcontrolador nRF52840 e se comunicam pelo protocolo
Enhanced ShockBurst a 20\,Hz; a cadeira funde as duas fontes, decide o alerta e
o devolve à unidade vestível, que aciona um motor de vibração. A caracterização
em bancada mostrou que a recolocação do sensor desloca a referência de postura
ereta em até cerca de 22\degree{} e que os FSR402 se prestam a indicar contato e
pressão relativa, não força calibrada, o que levou à calibração individual por
botão e à normalização das leituras pela linha de base do usuário. A regra de
alerta usa o desvio angular em relação à referência e, pelos sensores do
assento, distingue reclinar com apoio de escorregar o quadril e detecta a
cadeira vazia. Em capturas rotuladas com oito participantes, avaliadas deixando
cada participante fora do ajuste dos parâmetros, o sistema alertou em 58 de 62
segmentos de postura prejudicial e em 92\,\% do tempo dessas posturas, com
alerta em 3\,\% do tempo em postura ereta ou reclinada, contra 29\,\% da regra
baseada apenas no IMU. A validação em hardware com motor ativo confirmou o
funcionamento integrado e revelou interferência da vibração na leitura
inercial, posteriormente corrigida em firmware e verificada apenas por
reprodução dos dados. Em demonstração pública com
dezenas de pessoas, a postura prejudicial foi detectada em praticamente todos
os casos observados. Os resultados indicam a viabilidade técnica da fusão de
pressão no assento e orientação do tronco para detecção postural, com
limitações quanto ao tamanho da amostra, à projeção isolada da cabeça e ao
consumo de energia, ainda não medido.

\vspace{0.8em}
\noindent\textbf{Palavras-chave:} Monitoramento postural; Ergonomia; Distúrbios
osteomusculares relacionados ao trabalho; Sensores de força resistivos;
Unidade de medição inercial; Fusão de sensores.

\clearpage
\section*{Abstract}
\noindent
Work-related musculoskeletal disorders affect most office workers and are
associated with poor postures held for long periods. This work presents an
autonomous system that detects harmful sitting posture and alerts the user in
real time through vibration. It combines an office chair instrumented with
eight force-sensing resistors (FSR402), four on the seat and four on the
backrest, and a wearable unit with an inertial measurement unit (IMU) placed
just below the neck. Both run on nRF52840 microcontrollers and communicate over
Enhanced ShockBurst at 20\,Hz; the chair fuses both sources, decides the alert
and sends it back to the wearable, which drives a vibration motor. Bench
characterization showed that re-mounting the sensor shifts the upright
reference by up to about 22\degree{} and that the FSR402 is suited to contact
and relative pressure rather than calibrated force, which led to a per-user,
button-triggered calibration and to normalizing the readings by the user's
baseline. The alert rule uses the angular deviation from the reference and,
through the seat sensors, tells supported reclining from a forward slide of the
hips and detects an empty chair. On labeled captures with eight participants,
evaluated leaving each participant out of parameter tuning, the system alerted
in 58 of 62 harmful-posture segments and during 92\,\% of their time, while
alerting during 3\,\% of upright or reclined time, against 29\,\% for the
IMU-only rule. Hardware validation with the motor active confirmed end-to-end
operation and revealed vibration interference on the inertial reading, later
fixed in firmware and so far checked only by replaying the data. In a public demonstration with dozens of people, harmful
posture was detected in practically every observed case. The results indicate
the technical feasibility of fusing seat pressure and trunk orientation for
posture detection, with limitations regarding sample size, isolated head
protraction and the still unmeasured power consumption.

\vspace{0.8em}
\noindent\textbf{Keywords:} Posture monitoring; Ergonomics; Work-related
musculoskeletal disorders; Force-sensing resistors; Inertial measurement unit;
Sensor fusion.
